% GENERATED FILE. Edit canonical lessons, then run npm run generate:books.
% canonical-source-hash: fnv1a64:f6bb83adf5db0616
% canonical-lessons: ML-C409-nisedhikkuka, ML-C409-avakasappetuka, ML-C409-aropikkuka, ML-C409-nispaksamaya, ML-C409-visadikaranam

\chapter{Blame and Denial}
\label{ch:ml-blame-and-denial}
\hlchaptermodality{409}

\begin{chapteropening}
\textbf{By the end of this chapter:} \emph{I can say to deny, to claim, to allege, impartial and an explanation.}

Complete the last lesson of chapter 409: I can say to deny, to claim, to allege, impartial and an explanation.
\end{chapteropening}

\section[\texorpdfstring{\ml{നിഷേധിക്കുക} (niṣēdhikkuka)}{നിഷേധിക്കുക (niṣēdhikkuka)}]{\ml{നിഷേധിക്കുക} (niṣēdhikkuka) — to deny}

\label{lesson:ML-C409-nisedhikkuka}



\begin{quote}
Before the new one: say the Malayalam for an attitude, then the Malayalam for to justify.
\end{quote}

\subsection*{You'll want to know: \ml{നിഷേധിക്കുക}}
\textbf{\ml{നിഷേധിക്കുക}} — \emph{niṣēdhikkuka} — \textquotedblleft{}to deny\textquotedblright{}.

\begin{quote}
\ml{അവൻ} \ml{ആ} \ml{വാർത്ത} \ml{നിഷേധിച്ചു}.
\end{quote}

\begin{quote}
\emph{avan ā vārtta niṣēdhiccu.}
\end{quote}

\begin{quote}
He denied that news report.
\end{quote}

Past: \ml{നിഷേധിച്ചു}, \emph{niṣēdhiccu}.

The noun is \ml{നിഷേധം}, a denial; \ml{നിഷേധിക്കുക} can also mean to refuse a thing to someone.

\begin{cousinweb}
\textbf{\ml{നിഷേധിക്കുക}} (\emph{niṣēdhikkuka}) is built from Sanskrit \emph{niṣedha}, \textquotedblleft{}warding off, prevention, prohibition\textquotedblright{} (Monier-Williams), with the Malayalam verb ending -\ml{ിക്കുക}. To deny is to push a claim away.
\end{cousinweb}

\begin{grammarlens}[title={What you've built}]
One more word for giving reasons and opinions.
\end{grammarlens}

\subsection*{Your turn}
\begin{itemize}
\raggedright
  \item \emph{Say it:} \emph{niṣēdhikkuka}
  \item \emph{Say it:} \emph{niṣēdhikkuka}, once more
  \item \emph{Say it:} \emph{avan ā vārtta niṣēdhiccu.} — He denied that news report.
  \item \emph{Recall:} say the Malayalam for an attitude, then the Malayalam for to justify, then say \emph{niṣēdhikkuka} again
\end{itemize}

\begin{tcolorbox}[breakable,colback=teal!4,colframe=teal!35!black,title={Before you move on}]
What does \ml{നിഷേധിക്കുക} mean? (\textquotedblleft{}To deny\textquotedblright{}.) Say it once more.
\end{tcolorbox}
\section[\texorpdfstring{\ml{അവകാശപ്പെടുക} (avakāśappeṭuka)}{അവകാശപ്പെടുക (avakāśappeṭuka)}]{\ml{അവകാശപ്പെടുക} (avakāśappeṭuka) — to claim}

\label{lesson:ML-C409-avakasappetuka}



\begin{quote}
Before the new one: say the Malayalam for to justify, then the Malayalam for to deny.
\end{quote}

\subsection*{You'll want to know: \ml{അവകാശപ്പെടുക}}
\textbf{\ml{അവകാശപ്പെടുക}} — \emph{avakāśappeṭuka} — \textquotedblleft{}to claim\textquotedblright{}.

\begin{quote}
\ml{ആ} \ml{ഭൂമി} \ml{തന്റേതാണെന്ന്} \ml{അയാൾ} \ml{അവകാശപ്പെട്ടു}.
\end{quote}

\begin{quote}
\emph{ā bhūmi tanṟētāṇennŭ ayāḷ avakāśappeṭṭu.}
\end{quote}

\begin{quote}
He claimed that the land was his.
\end{quote}

Past: \ml{അവകാശപ്പെട്ടു}, \emph{avakāśappeṭṭu}.

Built on \ml{അവകാശം}, a right: to claim is to say a right is yours. What is claimed comes before \ml{എന്ന്}.

\begin{cousinweb}
\textbf{\ml{അവകാശപ്പെടുക}} (\emph{avakāśappeṭuka}) is built on \ml{അവകാശം}, \textquotedblleft{}a right\textquotedblright{}, from Sanskrit \emph{avakāśa}, \textquotedblleft{}space, room; opportunity\textquotedblright{} (CDIAL 728), with -\ml{പ്പെടുക} from native \ml{പെടുക} (DEDR 3853).
\end{cousinweb}

\begin{grammarlens}[title={What you've built}]
One more word for giving reasons and opinions.
\end{grammarlens}

\subsection*{Your turn}
\begin{itemize}
\raggedright
  \item \emph{Say it:} \emph{avakāśappeṭuka}
  \item \emph{Say it:} \emph{avakāśappeṭuka}, once more
  \item \emph{Say it:} \emph{ā bhūmi tanṟētāṇennŭ ayāḷ avakāśappeṭṭu.} — He claimed that the land was his.
  \item \emph{Recall:} say the Malayalam for to justify, then the Malayalam for to deny, then say \emph{avakāśappeṭuka} again
\end{itemize}

\begin{tcolorbox}[breakable,colback=teal!4,colframe=teal!35!black,title={Before you move on}]
What does \ml{അവകാശപ്പെടുക} mean? (\textquotedblleft{}To claim\textquotedblright{}.) Say it once more.
\end{tcolorbox}
\section[\texorpdfstring{\ml{ആരോപിക്കുക} (ārōpikkuka)}{ആരോപിക്കുക (ārōpikkuka)}]{\ml{ആരോപിക്കുക} (ārōpikkuka) — to allege; to accuse}

\label{lesson:ML-C409-aropikkuka}



\begin{quote}
Before the new one: say the Malayalam for to deny, then the Malayalam for to claim.
\end{quote}

\subsection*{You'll want to know: \ml{ആരോപിക്കുക}}
\textbf{\ml{ആരോപിക്കുക}} — \emph{ārōpikkuka} — \textquotedblleft{}to allege; to accuse\textquotedblright{}.

\begin{quote}
\ml{അയൽക്കാരൻ} \ml{പണം} \ml{എടുത്തു} \ml{എന്ന്} \ml{അവൾ} \ml{ആരോപിച്ചു}.
\end{quote}

\begin{quote}
\emph{ayalkkāran paṇaṁ eṭuttu ennŭ avaḷ ārōpiccu.}
\end{quote}

\begin{quote}
She alleged that the neighbour had taken the money.
\end{quote}

Past: \ml{ആരോപിച്ചു}, \emph{ārōpiccu}.

The noun is \ml{ആരോപണം}, an allegation. Against someone: -\ml{ന്} \ml{എതിരെ}.

\begin{cousinweb}
\textbf{\ml{ആരോപിക്കുക}} (\emph{ārōpikkuka}) is built from Sanskrit \emph{āropa}, \textquotedblleft{}laying on (as a burden), charging with\textquotedblright{} (Monier-Williams), with the Malayalam verb ending -\ml{ിക്കുക}: to allege is to lay a charge on someone.
\end{cousinweb}

\begin{grammarlens}[title={What you've built}]
One more word for giving reasons and opinions.
\end{grammarlens}

\subsection*{Your turn}
\begin{itemize}
\raggedright
  \item \emph{Say it:} \emph{ārōpikkuka}
  \item \emph{Say it:} \emph{ārōpikkuka}, once more
  \item \emph{Say it:} \emph{ayalkkāran paṇaṁ eṭuttu ennŭ avaḷ ārōpiccu.} — She alleged that the neighbour had taken the money.
  \item \emph{Recall:} say the Malayalam for to deny, then the Malayalam for to claim, then say \emph{ārōpikkuka} again
\end{itemize}

\begin{tcolorbox}[breakable,colback=teal!4,colframe=teal!35!black,title={Before you move on}]
What does \ml{ആരോപിക്കുക} mean? (\textquotedblleft{}To allege; to accuse\textquotedblright{}.) Say it once more.
\end{tcolorbox}
\section[\texorpdfstring{\ml{നിഷ്പക്ഷമായ} (niṣpakṣamāya)}{നിഷ്പക്ഷമായ (niṣpakṣamāya)}]{\ml{നിഷ്പക്ഷമായ} (niṣpakṣamāya) — impartial; neutral}

\label{lesson:ML-C409-nispaksamaya}



\begin{quote}
Before the new one: say the Malayalam for to claim, then the Malayalam for to allege.
\end{quote}

\subsection*{You'll want to know: \ml{നിഷ്പക്ഷമായ}}
\textbf{\ml{നിഷ്പക്ഷമായ}} — \emph{niṣpakṣamāya} — \textquotedblleft{}impartial; neutral\textquotedblright{}.

\begin{quote}
\ml{അന്വേഷണം} \ml{നിഷ്പക്ഷമായിരിക്കണം}.
\end{quote}

\begin{quote}
\emph{anvēṣaṇaṁ niṣpakṣamāyirikkaṇaṁ.}
\end{quote}

\begin{quote}
The inquiry must be impartial.
\end{quote}

The noun is \ml{നിഷ്പക്ഷത}, impartiality.

\begin{cousinweb}
\textbf{\ml{നിഷ്പക്ഷമായ}} (\emph{niṣpakṣamāya}) is built from Sanskrit \emph{niṣ-}, \textquotedblleft{}without\textquotedblright{}, and \emph{pakṣa}, \textquotedblleft{}a wing, a side\textquotedblright{} (CDIAL 7627), with the adjective ending -\ml{ആയ}: taking no side.
\end{cousinweb}

\begin{grammarlens}[title={What you've built}]
One more word for giving reasons and opinions.
\end{grammarlens}

\subsection*{Your turn}
\begin{itemize}
\raggedright
  \item \emph{Say it:} \emph{niṣpakṣamāya}
  \item \emph{Say it:} \emph{niṣpakṣamāya}, once more
  \item \emph{Say it:} \emph{anvēṣaṇaṁ niṣpakṣamāyirikkaṇaṁ.} — The inquiry must be impartial.
  \item \emph{Recall:} say the Malayalam for to claim, then the Malayalam for to allege, then say \emph{niṣpakṣamāya} again
\end{itemize}

\begin{tcolorbox}[breakable,colback=teal!4,colframe=teal!35!black,title={Before you move on}]
What does \ml{നിഷ്പക്ഷമായ} mean? (\textquotedblleft{}Impartial; neutral\textquotedblright{}.) Say it once more.
\end{tcolorbox}
\section[\texorpdfstring{\ml{വിശദീകരണം} (viśadīkaraṇaṁ)}{വിശദീകരണം (viśadīkaraṇaṁ)}]{\ml{വിശദീകരണം} (viśadīkaraṇaṁ) — an explanation}

\label{lesson:ML-C409-visadikaranam}



\begin{quote}
Before the new one: say the Malayalam for to allege, then the Malayalam for impartial.
\end{quote}

\subsection*{You'll want to know: \ml{വിശദീകരണം}}
\textbf{\ml{വിശദീകരണം}} — \emph{viśadīkaraṇaṁ} — \textquotedblleft{}an explanation\textquotedblright{}.

\begin{quote}
\ml{വൈകിയതിന്} \ml{അവൻ} \ml{ഒരു} \ml{വിശദീകരണവും} \ml{നൽകിയില്ല}.
\end{quote}

\begin{quote}
\emph{vaikiyatinŭ avan oru viśadīkaraṇavuṁ nalkiyilla.}
\end{quote}

\begin{quote}
He gave no explanation for being late.
\end{quote}

The verb \ml{വിശദീകരിക്കുക}, to explain, is the same word.

\begin{cousinweb}
\textbf{\ml{വിശദീകരണം}} (\emph{viśadīkaraṇaṁ}) is built from Sanskrit \emph{viśada}, \textquotedblleft{}clear, bright\textquotedblright{} (Monier-Williams), with \emph{-ī-karaṇa}, \textquotedblleft{}a making into\textquotedblright{}: a making clear. It is the noun of \ml{വിശദീകരിക്കുക}.
\end{cousinweb}

\begin{grammarlens}[title={What you've built}]
One more word for giving reasons and opinions.
\end{grammarlens}

\subsection*{Your turn}
\begin{itemize}
\raggedright
  \item \emph{Say it:} \emph{viśadīkaraṇaṁ}
  \item \emph{Say it:} \emph{viśadīkaraṇaṁ}, once more
  \item \emph{Say it:} \emph{vaikiyatinŭ avan oru viśadīkaraṇavuṁ nalkiyilla.} — He gave no explanation for being late.
  \item \emph{Recall:} say the Malayalam for to allege, then the Malayalam for impartial, then say \emph{viśadīkaraṇaṁ} again
\end{itemize}

\begin{tcolorbox}[breakable,colback=teal!4,colframe=teal!35!black,title={Before you move on}]
What does \ml{വിശദീകരണം} mean? (\textquotedblleft{}An explanation\textquotedblright{}.) Say it once more.
\end{tcolorbox}
